\subsection{Experimental result reproducibility}
\label{app:checklist_experimental_reproducibility}

\paragraph{Question.}
Are the main experimental results reproducible from the released paper, code,
configuration surfaces, scripts, generated benchmark definitions, random seeds,
and output artifacts?

\paragraph{Answer.}
Yes.

\paragraph{Justification.}
The paper-facing protocol specifies the benchmark populations, model families,
training budgets, evaluation horizons, statistical units, and seed aggregation
rules. The repository also contains executable configuration surfaces that make
the protocol reproducible: \texttt{pyproject.toml} and \texttt{uv.lock} define
the environment, \texttt{tools/train.py} is the single training entry point,
\texttt{skae/config.py} defines model and optimizer defaults, and the paper
launchers write task tables whose rows become the exact command-line arguments
passed to \texttt{tools/train.py}. Each training run persists its resolved
\texttt{config.json}, checkpoints, metrics, and evaluation outputs, so a
reproduction can audit the actual hyperparameters used by any displayed row.

The only qualification is that exact bitwise equality is not promised across
different GPU models, CUDA libraries, or PyTorch kernels. Reproducibility should
therefore be interpreted as reproducing the same task table, seeds, data
generation protocol, checkpoint-selection rule, and statistical summaries, with
small numerical differences possible across hardware/software stacks.

\paragraph{Environment.}
Use the locked project environment:
\begin{itemize}
  \item Python: \(\geq 3.10\), managed with \texttt{uv}.
  \item Dependency source: \texttt{pyproject.toml} plus \texttt{uv.lock}.
  \item Core dependencies: \texttt{torch>=2.2}, \texttt{numpy},
  \texttt{scipy}, \texttt{matplotlib}, \texttt{dysts>=0.96},
  \texttt{numba}, and \texttt{scikit-learn}. The Dysts appendix records the
  paper-facing resolved Dysts version as \texttt{0.96}.
  \item GPU jobs: the benchmark array runner uses SLURM on the \texttt{long}
  partition with one GPU, four CPUs, 16 GB memory, and \texttt{cuda/12.6.0}.
  CPU reducers and collectors are launched as SLURM jobs as well.
\end{itemize}

The setup command is:
\begin{verbatim}
uv sync
\end{verbatim}

All Python entry points should be run through \texttt{uv run}; paper-scale
training and evaluation should be submitted through the SLURM scripts rather
than run on a login node.

\paragraph{Core training and evaluation surfaces.}
The central command constructed by the launchers is:
\begin{verbatim}
uv run python tools/train.py \
  --config <config_name> \
  --env <env_name> \
  --env_dt <stored_step> \
  --num_steps <steps> \
  --batch_size 256 \
  --target_size 256 \
  --res_coeff 1.0 \
  --reconst_coeff 0.03 \
  --pred_coeff 1.0 \
  --sparsity_coeff <lambda_sp> \
  --sequence_length <L> \
  --eval_profile full \
  --seed <seed> \
  --device cuda \
  --log_dir <run_root>
\end{verbatim}

The model rows are selected by \texttt{--config}, \texttt{--k\_structure},
\texttt{--soft\_block}, LISTA flags, and the sparsity coefficient:
\begin{itemize}
  \item LISTA: \texttt{lista\_parity\_generic\_sparse} with
  \texttt{lista\_alpha=0.15}, \texttt{lista\_num\_loops=2},
  \texttt{lista\_final\_op=sign\_split}, dense \(K\), and
  \(\lambda_{\rm sp}=0.003\) on the multibasin benchmark.
  \item LISTA-BD: the same LISTA encoder, with
  \texttt{k\_structure=block\_diagonal}. The block count is set from benchmark
  metadata only for this architecture diagnostic.
  \item LISTA-SB: the same LISTA encoder, dense \(K\), and a soft off-block
  penalty. The soft-block count is likewise architecture metadata, not a
  training label.
  \item Sparse MLP: \texttt{generic\_sparse}, dense \(K\), and the paper
  sparsity coefficient.
  \item Sparse MLP-BD: \texttt{generic\_sparse} with block-diagonal \(K\).
  \item Dense MLP: \texttt{generic\_no\_shrink} with \(\lambda_{\rm sp}=0\)
  and no final ReLU sparsifying gate.
\end{itemize}

\paragraph{Multibasin benchmark route.}
The paper-facing controlled benchmark contains 15 two-dimensional multibasin
systems:
\begin{verbatim}
gated_local_linear
gated_transfer_linear
claude:arrested_spiral
claude:cal_asymmetric_3
claude:cal_high_cross_3
claude:cal_hexagon_6
claude:cal_octagon_8
claude:cal_pentagon_5
claude:cal_square_4
claude:duffing_triple_well
claude:snic_multi
claude:transition_routes_4
claude:var_depth_gradient_4
claude:var_diamond_4
claude:var_l_shape_5
\end{verbatim}

They are continuous-time vector fields integrated with fourth-order Runge--Kutta
to produce stored observations. Training observes only stored states; vector
fields, integration substeps, basin labels, basin counts, and trajectory-to-basin
assignments are not supplied to the encoder or transition model. Basin metadata
is used after training for evaluation and, only for the block-diagonal and
soft-block diagnostic rows, to set the fixed number of transition groups.

The multibasin paper settings are:
\begin{itemize}
  \item Seeds: \(0,1,\ldots,14\).
  \item Latent dimension: \(d_z=256\).
  \item Rollout-window length: \(L=8\), i.e. nine adjacent stored states.
  \item Optimization: \(200{,}000\) steps, batch size \(256\), AdamW, encoder
  and decoder learning rate \(5\times10^{-5}\), Koopman-matrix learning rate
  \(5\times10^{-6}\), and weight decay \(10^{-4}\) on non-transition
  parameters.
  \item Loss weights: \(\lambda_{\rm pred}=1\), \(\lambda_{\rm rec}=0.03\),
  \(\lambda_{\rm lin}=1\), \(\lambda_{\rm sp}=0.003\) for sparse rows, and
  \(\lambda_{\rm sp}=0\) for Dense MLP.
  \item Training-state distribution: boundary-emphasized hard-initial-condition
  oversampling enabled for all six controlled rows; \(50\%\) of windows start
  from a retained pool of 1024 states selected from 4096 candidates using
  32-step, label-free perturbation/transient scores, with four perturbations,
  perturbation scale \(0.04\), transient window \(8\), transient weight \(0.5\),
  and jitter scale \(0.25\).
  \item Checkpoint rule: every 500 steps, and at the final step, evaluate a
  200-step every-step-reencoded validation rollout from 16 held-out initial
  states; save the checkpoint with lowest final validation error as
  \texttt{checkpoint.pt}.
\end{itemize}

The reproducibility route is to regenerate task tables and submit them through
the existing array runner:
\begin{verbatim}
uv run python tools/build_transition_rich_basin_partition_tasks.py \
  --phase_label transition_rich_basin_partition \
  --output_tsv <task_tsv> \
  --output_manifest_json <manifest_json> \
  --systems_csv <the 15-system comma-separated list above> \
  --model_variants_csv lista_blockdiag_signsplit_hardinit_basin_partition,\
mlp_sparse_blockdiag_hardinit_basin_partition_control,\
mlp_sparse_hardinit_basin_partition_control,\
mlp_zero_sparse_hardinit_basin_partition_control \
  --seeds_csv 0,1,2,3,4,5,6,7,8,9,10,11,12,13,14 \
  --num_steps_override 200000

TASK_TSV=<task_tsv> BASE_OUT=<scratch_output_root> \
  sbatch --array=0-<num_tasks_minus_1> scripts/run_paper_benchmark_array.sh
\end{verbatim}

The current repository also includes paper-facing convenience launchers for the
matched-dimension LISTA, LISTA-SB, and backfill rows:
\begin{verbatim}
sbatch scripts/queue_transition_rich_lista_sb_p256_hardinit_fairness.sh
sbatch scripts/queue_transition_rich_lista_dense_p256_hardinit_table123.sh
sbatch scripts/queue_transition_rich_table2_5model_seed_backfill.sh
sbatch scripts/queue_sparse_mlp_bd_repaired_table1.sh
bash scripts/queue_per_basin_deep_current_table1_eval.sh
\end{verbatim}

Some older or broader scripts still know about 17-system exploratory packets.
For reproducing the claims in this paper, use the 15-system roster above or the
filtered paper-facing figure/statistics scripts.

\paragraph{Dysts \(dt\times30\) forecasting route.}
The paper-facing Dysts benchmark contains these 10 three-dimensional systems:
\begin{verbatim}
dysts:Chua
dysts:Dadras
dysts:DequanLi
dysts:Hadley
dysts:LuChenCheng
dysts:QiChen
dysts:Sakarya
dysts:SanUmSrisuchinwong
dysts:ShimizuMorioka
dysts:WangSun
\end{verbatim}

For each system, stored observations use \(30\) times the native Dysts
integration interval, and coordinates are standardized after trajectory
generation. The full Dysts cache profile uses 200 cached trajectories, 30000
stored observations per trajectory, and a 2000-step warm-up. Initial conditions
are the Dysts default initial condition plus Gaussian perturbations with scale
\(0.2\) times the coordinate-wise Dysts standard deviation. Train, validation,
and test caches are split by cache namespace; held-out test windows are disjoint
from training windows.

The Dysts paper settings are:
\begin{itemize}
  \item Seeds: \(0,1,\ldots,14\).
  \item Latent dimension: \(d_z=256\).
  \item Rollout-window length: \(L=10\).
  \item Optimization: \(100{,}000\) steps and batch size \(256\).
  \item LISTA learning rates: \(5\times10^{-5}\) for encoder/decoder and
  \(5\times10^{-6}\) for \(K\).
  \item MLP learning rates: \(10^{-4}\) for encoder/decoder and \(10^{-5}\)
  for \(K\).
  \item Sparse-row Dysts coefficient: \(\lambda_{\rm sp}=0.006\); Dense MLP:
  \(\lambda_{\rm sp}=0\).
  \item Long-horizon evaluation: 100 held-out test windows per system and seed,
  horizons \(100,500,1000,1500,2000,3000,4000,5000\), and periodic
  re-encoding periods \(10,25,50,100,150,200\).
\end{itemize}

Use the Dysts queue with the 10-system paper list explicitly, because the task
builder also contains additional stress-test systems:
\begin{verbatim}
SYSTEMS_CSV=dysts:Chua,dysts:Dadras,dysts:DequanLi,dysts:Hadley,\
dysts:LuChenCheng,dysts:QiChen,dysts:Sakarya,\
dysts:SanUmSrisuchinwong,dysts:ShimizuMorioka,dysts:WangSun \
SEEDS_CSV=0,1,2,3,4,5,6,7,8,9,10,11,12,13,14 \
NUM_STEPS=100000 \
sbatch scripts/queue_dysts_dt30_basinblock_p256_seeds0to14.sh
\end{verbatim}

That launcher builds the task table with
\texttt{tools/build\_dysts\_dt30\_basinblock\_tasks.py}, prebuilds Dysts
\texttt{train}/\texttt{val} caches, runs packed training arrays through
\texttt{scripts/run\_paper\_benchmark\_packed\_array.sh}, and queues
\texttt{scripts/queue\_dysts\_long\_horizon\_eval.sh} for the held-out test
cache evaluation.

\paragraph{Support and mechanism diagnostics.}
Support definitions are fixed in code and prose:
\begin{itemize}
  \item \(S_{\rm abs}\): active coordinates with absolute threshold \(10^{-3}\).
  \item \(F_{\rm abs}\): greedy Jaccard families of \(S_{\rm abs}\), with
  threshold \(0.5\), used for the main basin-support alignment diagnostic.
  \item \(S_{\rm top8}\) and \(F_{\rm top8}\): eight largest-magnitude
  coordinates and their support families, used for routing and refresh
  appendix diagnostics.
\end{itemize}

The per-basin deep-state slice is evaluation-only: within each benchmark basin,
states are ranked by the basin-depth margin, defined as the distance to the
second-nearest attractor center minus the distance to the nearest attractor
center, and the top quartile is retained. Support reducers use
\texttt{eval\_seed=42} by default, \(128\) trajectories of length \(128\),
endpoint rollouts of \(5000\) steps for basin labeling when needed, and
\texttt{DEPTH\_SLICE\_MODE=per\_basin} for the current paper table.

The coordinate-intervention experiment uses the LISTA dense-transition model,
the \texttt{gated\_local\_linear} system, training seed 0, 100 held-out
basin-interior starts, horizons \(1,3,\ldots,21\), top-\(k\) coordinate drops
up to \(k=10\), and 20 random inactive-coordinate reassignments. The route is:
\begin{verbatim}
NUM_INITIAL_POINTS=100 \
RANDOM_SUPPORT_REPEATS=20 \
ROWS_CSV=<paper_forecasting_rows.csv> \
OUT_DIR=<support_intervention_output_dir> \
sbatch scripts/run_support_coordinate_interventions.sh
\end{verbatim}

\paragraph{Randomness and seeds.}
Training seeds are the explicit model seeds \(0,\ldots,14\). Within a run,
\texttt{tools/train.py} sets \texttt{torch.manual\_seed(cfg.SEED)} and, when
available, the CUDA seed. Online training batches use deterministic
\texttt{torch.Generator} instances derived from the run seed. The fixed
validation set uses seed \(\texttt{cfg.SEED}+999999\). Standardized evaluation
uses seed \(\texttt{cfg.SEED}+12345\). Dysts cache generation is intentionally
independent of the training seed and uses split-specific deterministic cache
seeds: train, validation, and test are separate namespaces. Support reducers use
default evaluation seed \(42\), and random-support shuffling uses random seed
\(123\) unless overridden.

\paragraph{Outputs and paper regeneration.}
Each run directory contains:
\begin{itemize}
  \item \texttt{config.json}: resolved configuration and command-line overrides.
  \item \texttt{checkpoint.pt}: best checkpoint by the validation rule.
  \item \texttt{last.pt}: final checkpoint.
  \item \texttt{metrics\_history.jsonl}, \texttt{metrics\_summary.json},
  \texttt{final\_metrics.json}, and \texttt{training\_metrics.png}.
  \item \texttt{evaluation\_results\_best.json},
  \texttt{evaluation\_results\_last.json}, \texttt{evaluation\_summary.json},
  and per-system evaluation directories with optional rollout artifacts.
\end{itemize}

Collectors and reducers then produce paper-facing tables:
\begin{itemize}
  \item Forecasting rows: \texttt{forecasting\_rows.csv}, produced by
  \texttt{tools/collect\_forecasting\_roots.py} through
  \texttt{scripts/collect\_transition\_rich\_basin\_partition.sh},
  \texttt{scripts/collect\_paper\_benchmark.sh}, or the Dysts long-horizon
  collector.
  \item Support diagnostics: \texttt{interpretability\_rows.csv}, produced by
  \texttt{scripts/reduce\_transition\_rich\_interpretability\_metrics.sh} and
  merged by \texttt{scripts/merge\_transition\_rich\_interpretability\_shards.sh}.
  \item Dysts long-horizon rows and compact rollout artifacts from
  \texttt{tools/evaluate\_dysts\_long\_horizon\_run.py} and
  \texttt{scripts/collect\_dysts\_long\_horizon\_forecasting.sh}.
  \item Final paper figures and LaTeX tables under
  \texttt{docs/figures/neurips\_paper\_2026/} and
  \texttt{docs/figures/neurips\_paper\_2026/\_tables/}, generated by
  \texttt{scripts/build\_per\_system\_stats\_and\_forest.py} and
  \texttt{scripts/build\_paper\_seed\_stats\_and\_figures.py}.
\end{itemize}

Point estimates are interquartile means over seeds within each system followed
by arithmetic means across systems, except for explicitly descriptive count
columns. Multibasin forecasting uses paired within-system Wilcoxon signed-rank
tests with Holm correction. Dysts forecasting uses per-system seed-IQM effects
and a one-sided exact sign test across the 10 systems with Holm correction.
